% !PRINTED 18
\weberpage{18}

Wenn man also $\mu = \frac{1}{12}$ setzt, so gehen die Gleichungen (19)
in die Identitäten $g_2=g_2$, $g_3=g_3$ über.

Wenn man dann $y=1$ setzt, so ergibt die Transformation (18)

\weq{20}{\dfrac{2\,dx}{\sqrt{4\,x^3-g_2\,x-g_3}} = \dfrac{dz}{\sqrt{4\,z^3-g_2\,z-g_3}}\,.}

% [sic] Printed "nach (16)"; with \mu=1/12, the displayed substitution follows from (17).
Diese Transformation wird nach (16) durch

\[ z = \dfrac{-H}{48\,f} \]

vermittelt. Es ergibt sich aber aus (14):

\[ H = -48\left[(x^2+\tfrac{1}{4}g_2)^2 + 2\,g_3\,x\right], \]

und folglich erhalten wir

\weq{21}{z = \dfrac{(x^2+\tfrac{1}{4}g_2)^2+2\,g_3\,x}{4\,x^3-g_2\,x-g_3}\,.}

Weiter ergibt sich noch aus (14)

\[ T = 64\,x^6 - 80\,g_2\,x^4 - 320\,g_3\,x^3 - 20\,g_2^2\,x^2 - 16\,g_2\,g_3\,x + g_2^3 - 32\,g_3^2 \]

und dann nach (16)

\weq{22}{4\sqrt{4\,x^3-g_2\,x-g_3}^{\,3}\ \sqrt{4\,z^3-g_2\,z-g_3} = T\,.}

\gesperrt{Durch (20), (21), (22) ist die Multiplikation des elliptischen
 Differentials mit 2 geleistet.} Es ist dadurch nicht nur
$z$ als eindeutige (rationale) Funktion von $x$ dargestellt,
sondern auch die eine Quadratwurzel eindeutig durch die andere,
d.\,h. es ist einem \gesperrt{Punkte} $x$ eindeutig ein \gesperrt{Punkt} $z$ zugeordnet.

\begin{center}
§\,6. \gesperrt{Elliptische Kurven.}
\end{center}

Jede algebraische Abhängigkeit zwischen zwei Veränderlichen
$x$, $y$ wird ausgedrückt durch eine Gleichung der Form

\weq{1}{F(x,y) = 0\,,}

worin $F(x,y)$ eine ganze Funktion der beiden Veränderlichen $x$, $y$
bezeichnet. Betrachtet man $x$ und $y$ als Cartesische Koordinaten
eines Punktes in der Ebene, so ist (1) die Gleichung einer Kurve
$n$ten Grades, wenn $F$ in bezug auf $x$ und $y$ zusammengenommen
von der $n$ten Dimension ist. Diese Kurve ist dann das geometrische
 Bild der algebraischen Abhängigkeit, wobei indessen auch
imaginäre Punkte mit berücksichtigt werden müssen.
